\documentclass[12pt]{article}

\usepackage[margin=1in]{geometry}
\usepackage{enumitem}
\usepackage{amsmath}
\usepackage{amssymb}
\usepackage{array}
\usepackage{titlesec}

\title{Lecture 6 Notes:\\ Aristotle's \textit{Posterior Analytics}, Book I}
\author{}
\date{}

\begin{document}

\maketitle

\section*{Central Theme}

Aristotle's \textit{Posterior Analytics}, Book I is the transition from valid
inference to scientific knowledge.

In the \textit{Prior Analytics}, Aristotle asked:

\[
\text{When do propositions force a conclusion?}
\]

In the \textit{Posterior Analytics}, he asks:

\[
\text{When does a valid conclusion count as knowledge?}
\]

\[
\boxed{
\textit{Posterior Analytics}, \text{ Book I}
=
\text{Aristotle's theory of scientific demonstration}
}
\]

A syllogism gives a conclusion. A demonstration gives knowledge of the reason
why.

\section*{1. From Syllogism to Demonstration}

The \textit{Prior Analytics} studied syllogism in general. The
\textit{Posterior Analytics} studies a special kind of syllogism:
demonstration.

\[
\textit{Prior Analytics} = \text{valid syllogism}
\]

\[
\textit{Posterior Analytics} = \text{scientific demonstration}
\]

A demonstration is a syllogism productive of scientific knowledge.

\[
\boxed{
\text{Demonstration} =
\text{a syllogism that produces scientific knowledge}
}
\]

\subsection*{Teaching Point}

Not every valid syllogism is a demonstration.

A syllogism may be formally valid and still fail to produce scientific
knowledge if its premises are not true, explanatory, primary, necessary, and
appropriate.

\[
\boxed{
\text{Validity} \neq \text{demonstration}
}
\]

\section*{2. All Teaching Begins from Pre-Existent Knowledge}

Aristotle begins with the claim that all instruction given or received by
argument proceeds from pre-existent knowledge.

This applies to:

\begin{itemize}[itemsep=0.25em]
    \item mathematics,
    \item speculative sciences,
    \item syllogistic reasoning,
    \item induction,
    \item rhetoric,
    \item example,
    \item enthymeme.
\end{itemize}

\subsection*{Two Kinds of Pre-Existent Knowledge}

Sometimes the learner must already know that something is.

\[
\text{knowledge that something exists}
\]

Sometimes the learner must already know what a term means.

\[
\text{knowledge of meaning}
\]

Sometimes both are required.

\subsection*{Examples}

With triangle, one may need to understand what the word means.

\[
\text{triangle} = \text{a three-sided rectilinear figure}
\]

With unit, one must understand the meaning and also assume that such a thing
exists.

\[
\text{unit} = \text{meaning} + \text{existence}
\]

\subsection*{Teaching Point}

Learning never begins from absolutely nothing.

\[
\boxed{
\text{Before learning, we must already know something.}
}
\]

\section*{3. Aristotle and the Puzzle of the \textit{Meno}}

Aristotle's opening claim helps answer the problem made famous in Plato's
\textit{Meno}.

The puzzle is:

\begin{quote}
If we already know, we do not need to learn. If we do not know, we cannot know
what to look for.
\end{quote}

Aristotle answers by distinguishing different senses of knowing.

A person may know something universally without yet recognizing its application
to a particular case.

\subsection*{Example}

A student may know:

\[
\text{Every triangle has angles equal to two right angles.}
\]

But the student may not yet know that this figure here is a triangle.

Once the student recognizes the particular as falling under the universal, the
student comes to know the particular case.

\subsection*{Teaching Point}

Learning moves from prior universal grasp to particular recognition.

\[
\boxed{
\text{One may know universally and still learn the particular application.}
}
\]

\section*{4. What Is Scientific Knowledge?}

Aristotle says that we possess unqualified scientific knowledge when we know:

\begin{enumerate}[itemsep=0.25em]
    \item the cause on which the fact depends,
    \item that it is the cause of that fact,
    \item that the fact could not be otherwise.
\end{enumerate}

\[
\boxed{
\text{Scientific knowledge} =
\text{knowledge of the cause as cause}
}
\]

Scientific knowledge is not merely true belief. It is knowledge of why the
truth must be so.

\[
\boxed{
\text{Scientific knowledge concerns what cannot be otherwise.}
}
\]

\subsection*{Teaching Point}

To know scientifically is not simply to know that something is true. It is to
know why it is necessarily true.

\[
\text{knowledge that} \quad \neq \quad \text{knowledge why}
\]

\section*{5. Demonstration as Scientific Syllogism}

Demonstration is the kind of syllogism that produces scientific knowledge.

\[
\boxed{
\text{Demonstration} =
\text{syllogism productive of scientific knowledge}
}
\]

The premises of a demonstration must be:

\[
\begin{array}{ll}
1. & \text{true} \\
2. & \text{primary} \\
3. & \text{immediate} \\
4. & \text{better known than the conclusion} \\
5. & \text{prior to the conclusion} \\
6. & \text{causes of the conclusion}
\end{array}
\]

\subsection*{Teaching Point}

A demonstration is not just a valid argument. It is a valid argument whose
premises explain the conclusion.

\[
\boxed{
\text{Syllogism} + \text{explanatory premises}
=
\text{demonstration}
}
\]

\section*{6. Example: Validity Without Demonstration}

Consider the following argument:

\[
\text{All stones are animals.}
\]

\[
\text{All humans are stones.}
\]

\[
\therefore \text{All humans are animals.}
\]

The conclusion is true, and the form resembles a valid syllogism. But this is
not a demonstration.

Why not?

\begin{itemize}[itemsep=0.25em]
    \item The premises are false.
    \item The premises are not explanatory.
    \item The middle term does not give the cause.
    \item The argument does not show why humans are animals.
\end{itemize}

\subsection*{Teaching Point}

Demonstration requires more than a conclusion that happens to be true.

\[
\boxed{
\text{A true conclusion from bad premises is not scientific knowledge.}
}
\]

\section*{7. Prior and Better Known: To Us and By Nature}

Aristotle distinguishes two senses of being prior and better known.

\[
\text{prior and better known to us}
\]

\[
\text{prior and better known by nature}
\]

Things closer to sense-perception are usually better known to us.

\[
\text{particulars, examples, visible cases}
\]

Causes and universals are prior and better known by nature.

\[
\text{principles, causes, universals}
\]

\subsection*{Whiteboard}

\[
\begin{array}{ll}
\textbf{To us} & \text{sense-near particulars} \\
\textbf{By nature} & \text{causes and universals}
\end{array}
\]

\subsection*{Teaching Point}

Learning starts from what is clearer to us and aims at what is clearer by
nature.

\[
\boxed{
\text{Education moves from familiar facts to explanatory causes.}
}
\]

\section*{8. Basic Truths}

A basic truth is an immediate proposition.

\[
\boxed{
\text{Basic truth} = \text{immediate proposition}
}
\]

An immediate proposition has no prior proposition through which it is proved.

Basic truths are necessary because demonstration cannot go on forever. If every
premise required another premise before it, no demonstration could ever be
completed.

\subsection*{Teaching Point}

Demonstration rests on starting points that are not themselves demonstrated.

\[
\boxed{
\text{Science requires indemonstrable first principles.}
}
\]

\section*{9. Thesis, Axiom, Hypothesis, and Definition}

Aristotle distinguishes several kinds of starting points.

\[
\begin{array}{ll}
\textbf{Thesis} &
\text{an immediate basic truth not proved by the teacher} \\
\textbf{Axiom} &
\text{a truth the learner must know to learn anything} \\
\textbf{Hypothesis} &
\text{an assumption of existence or non-existence} \\
\textbf{Definition} &
\text{a statement of what something means}
\end{array}
\]

\subsection*{Hypothesis and Definition}

A hypothesis asserts that something is or is not.

\[
\text{There are units.}
\]

A definition states what something is without asserting its existence.

\[
\text{A unit is quantitatively indivisible.}
\]

\subsection*{Teaching Point}

A definition lays down meaning; a hypothesis asserts existence.

\[
\boxed{
\text{To define what something is is not yet to assert that it exists.}
}
\]

\section*{10. Against Infinite Regress}

Aristotle rejects the view that all knowledge requires demonstration.

If every premise had to be demonstrated by a prior premise, then knowledge would
require an infinite regress.

\[
P_1 \leftarrow P_2 \leftarrow P_3 \leftarrow P_4 \leftarrow \cdots
\]

But an infinite regress cannot be traversed.

\[
\boxed{
\text{Infinite regress} \Rightarrow \text{no completed scientific knowledge}
}
\]

Therefore, not all knowledge is demonstrative. Some knowledge must be knowledge
of immediate first principles.

\subsection*{Teaching Point}

Demonstration is possible only because some truths are known without being
demonstrated.

\[
\boxed{
\text{Not all knowledge is demonstrative.}
}
\]

\section*{11. Against Circular Demonstration}

Aristotle also rejects circular demonstration.

Circular demonstration attempts to prove premises by means of the conclusion
they are supposed to prove.

\[
A \Rightarrow B
\]

\[
B \Rightarrow A
\]

This destroys explanatory priority.

\subsection*{Why Circular Demonstration Fails}

Demonstration must proceed from what is prior and better known to what is
posterior and less known.

But in circular demonstration, the same thing becomes both prior and posterior.

\[
\boxed{
\text{The same thing cannot be both prior and posterior in the same demonstration.}
}
\]

\subsection*{Teaching Point}

Circular demonstration may show reciprocal implication, but it does not give
scientific explanation.

\[
\boxed{
\text{Science requires first principles, not explanatory circles.}
}
\]

\section*{12. Types of Attribute}

Aristotle next asks what kind of attributes can be demonstrated scientifically.

He distinguishes:

\[
\begin{array}{ll}
\textbf{True in every instance} & \text{true of all cases and always} \\
\textbf{Essential} & \text{belongs to what the subject is} \\
\textbf{Commensurately universal} & \text{belongs to the subject as such} \\
\textbf{Accidental} & \text{not essentially connected}
\end{array}
\]

\subsection*{Teaching Point}

Scientific demonstration is concerned with what belongs essentially and
universally, not with accidental conjunctions.

\[
\boxed{
\text{Demonstration proves essential attributes of their proper subjects.}
}
\]

\section*{13. True in Every Instance}

An attribute is true in every instance when it is predicable of all instances of
a subject and at all times.

\subsection*{Example}

\[
\text{Animal is predicable of every human.}
\]

If this is a human, then this is an animal.

\[
\text{Human} \Rightarrow \text{Animal}
\]

\subsection*{Teaching Point}

A scientific predicate cannot belong only occasionally or to only some cases.

\[
\boxed{
\text{Scientific demonstration requires stable predication.}
}
\]

\section*{14. Essential Attributes}

An attribute is essential in two especially important ways.

\subsection*{First Sense}

An attribute is essential when it belongs as an element in the subject's
essential nature.

\[
\text{Line belongs to triangle.}
\]

\[
\text{Point belongs to line.}
\]

\subsection*{Second Sense}

An attribute is essential when the subject is contained in the definition of the
attribute.

\[
\text{Odd belongs to number.}
\]

\[
\text{Even belongs to number.}
\]

The definition of odd or even includes number.

\subsection*{Teaching Point}

Essential attributes are not accidental add-ons. They are tied to what the
subject is.

\[
\boxed{
\text{What belongs essentially belongs through the thing's nature.}
}
\]

\section*{15. Accidental Attributes}

An accidental attribute is not connected with the subject through its essence.

\subsection*{Examples}

\[
\text{The animal is white.}
\]

\[
\text{The person is musical.}
\]

These may be true, but they do not follow from what animal or human is.

\subsection*{Teaching Point}

The accidental may be true, but it is not the proper object of scientific
demonstration.

\[
\boxed{
\text{There is no demonstration of the accidental as accidental.}
}
\]

\section*{16. Commensurately Universal Attributes}

An attribute belongs commensurately and universally when it belongs:

\begin{enumerate}[itemsep=0.25em]
    \item to every instance of the subject,
    \item essentially,
    \item as such,
    \item to the first proper subject.
\end{enumerate}

\[
\boxed{
\text{Commensurately universal} =
\text{universal, essential, and proper}
}
\]

\subsection*{Triangle Example}

Having angles equal to two right angles belongs to triangle as such.

It does not belong primarily to:

\[
\text{isosceles triangle}
\]

or:

\[
\text{scalene triangle}
\]

or:

\[
\text{equilateral triangle}
\]

It belongs first to:

\[
\text{triangle}
\]

\subsection*{Whiteboard}

\[
\text{isosceles} \rightarrow \text{triangle} \rightarrow \text{two right angles}
\]

\[
\boxed{
\text{The proper subject is the first subject to which the attribute belongs as such.}
}
\]

\section*{17. The Error of the Wrong Universal}

Aristotle warns that we may mistakenly think we have demonstrated something
commensurately and universally when we have not.

\subsection*{Example}

If we prove separately that:

\[
\text{equilateral triangles have angles equal to two right angles,}
\]

\[
\text{isosceles triangles have angles equal to two right angles,}
\]

\[
\text{scalene triangles have angles equal to two right angles,}
\]

we still have not grasped the attribute at the deepest scientific level unless
we know that it belongs to triangle as such.

\subsection*{Teaching Point}

Scientific knowledge seeks the correct level of universality.

\[
\boxed{
\text{Do not prove of a species what belongs first to the genus.}
}
\]

\section*{18. Necessary and Essential Premises}

Because scientific knowledge concerns what cannot be otherwise, demonstration
must proceed from necessary premises.

Essential attributes attach necessarily to their subjects. Accidental attributes
do not.

\[
\text{Essential} \Rightarrow \text{necessary}
\]

\[
\text{Accidental} \nRightarrow \text{necessary}
\]

\subsection*{Example}

\[
\text{A triangle has angles equal to two right angles.}
\]

This belongs necessarily.

\[
\text{This drawn triangle is blue.}
\]

This may be true, but not necessarily.

\subsection*{Teaching Point}

Demonstrative premises must be necessary and essential.

\[
\boxed{
\text{Science demonstrates necessary connections through essential premises.}
}
\]

\section*{19. Demonstration Must Remain Within One Genus}

Aristotle argues that demonstration cannot simply pass from one genus to
another.

For example, geometry cannot prove geometrical truths by arithmetic, unless the
sciences stand in a special subordinate relation.

\[
\boxed{
\text{Demonstration must use principles appropriate to its science.}
}
\]

\subsection*{The Three Elements of Demonstration}

Every demonstrative science involves:

\[
\begin{array}{ll}
1. & \text{the subject genus} \\
2. & \text{the axioms or basic premises} \\
3. & \text{the attributes demonstrated}
\end{array}
\]

\subsection*{Teaching Point}

Each science has its own domain, principles, and proper attributes.

\[
\boxed{
\text{Scientific knowledge is organized by genera.}
}
\]

\section*{20. Common Axioms and Proper Principles}

Some truths are common to many sciences, but they are used only within the
limits of a given science.

\subsection*{Example}

\[
\text{If equals are subtracted from equals, equals remain.}
\]

This may be common, but the geometer uses it of magnitudes, while the
arithmetician uses it of numbers.

\subsection*{Proper Principles}

Each science also has proper principles peculiar to its subject matter.

\[
\text{Geometry} \rightarrow \text{points, lines, figures}
\]

\[
\text{Arithmetic} \rightarrow \text{units, numbers}
\]

\subsection*{Teaching Point}

Common axioms do not erase the difference between sciences.

\[
\boxed{
\text{A science needs both common logical principles and proper subject principles.}
}
\]

\section*{21. Knowledge of the Fact and Knowledge of the Reasoned Fact}

Aristotle distinguishes:

\[
\text{knowledge of the fact}
\]

from:

\[
\text{knowledge of the reasoned fact}
\]

The first tells us that something is so. The second tells us why it is so.

\[
\boxed{
\text{Scientific knowledge, strictly speaking, is knowledge why.}
}
\]

\section*{22. The Planets Example}

Aristotle gives the example of the planets.

One might argue:

\[
\text{The planets do not twinkle.}
\]

\[
\text{Things that do not twinkle are near.}
\]

\[
\therefore \text{The planets are near.}
\]

This proves the fact, but not the reasoned fact.

The true explanatory order is:

\[
\text{The planets are near.}
\]

\[
\text{Near heavenly bodies do not twinkle.}
\]

\[
\therefore \text{The planets do not twinkle.}
\]

\subsection*{Teaching Point}

The cause must be the middle term.

\[
\boxed{
\text{In demonstration, the middle term is not merely a connector; it is the cause.}
}
\]

\section*{23. The Moon Example}

Another example concerns the moon.

One may infer:

\[
\text{The moon waxes in this way.}
\]

\[
\therefore \text{The moon is spherical.}
\]

But the deeper explanation is:

\[
\text{The moon waxes in this way because it is spherical.}
\]

\subsection*{Teaching Point}

The fact may be more obvious to us, while the cause is prior by nature.

\[
\boxed{
\text{The reasoned fact reverses the order of mere observation.}
}
\]

\section*{24. Subordinate Sciences}

Sometimes one science knows the fact while another knows the reason why.

\subsection*{Examples}

\[
\text{Optics} \rightarrow \text{subordinate to geometry}
\]

\[
\text{Harmonics} \rightarrow \text{subordinate to arithmetic}
\]

\[
\text{Observation} \rightarrow \text{subordinate to mathematical astronomy}
\]

The empirical observer may know that something is so, while the mathematical
science explains why.

\subsection*{Teaching Point}

Some sciences are related as fact-gathering to cause-explaining.

\[
\boxed{
\text{Subordinate sciences receive their explanations from superior sciences.}
}
\]

\section*{25. The First Figure Is the Scientific Figure}

Aristotle argues that the first figure is the most scientific.

Why?

\begin{enumerate}[itemsep=0.25em]
    \item It is the vehicle of mathematical demonstration.
    \item It best gives the cause.
    \item It enables knowledge of essence.
    \item The other figures depend on it.
\end{enumerate}

\[
\boxed{
\text{First figure} = \text{the figure of scientific explanation}
}
\]

\subsection*{Teaching Point}

In the \textit{Prior Analytics}, the first figure was the clearest form of
syllogism. In the \textit{Posterior Analytics}, it becomes the chief form of
scientific demonstration.

\[
\boxed{
\text{The most scientific syllogism is the one that best displays the cause.}
}
\]

\section*{26. Ignorance as Error}

Aristotle distinguishes ignorance as mere absence of knowledge from ignorance
as positive error produced by inference.

\[
\text{Ignorance as lack} = \text{not knowing}
\]

\[
\text{Ignorance as error} = \text{thinking falsely through inference}
\]

\subsection*{Teaching Point}

Bad reasoning can produce a positive state of ignorance.

\[
\boxed{
\text{Error is not merely absence; it can be false conviction.}
}
\]

\section*{27. Sense Perception and Induction}

Aristotle says that if a sense is lacking, a corresponding kind of knowledge is
also lacking.

Why?

Because induction depends on sense-perception, and universals are grasped
through induction.

\[
\text{sense perception} \rightarrow \text{particulars}
\]

\[
\text{particulars} \rightarrow \text{induction}
\]

\[
\text{induction} \rightarrow \text{universals}
\]

\[
\text{universals} \rightarrow \text{demonstration}
\]

\subsection*{Teaching Point}

There is no science without sense, but science is not mere sense.

\[
\boxed{
\text{Knowledge begins in perception and is completed in universal explanation.}
}
\]

\section*{28. The Regress of Premises Must Terminate}

A major question in Book I is whether demonstration can involve an infinite
regress of premises.

\[
A \rightarrow B \rightarrow C \rightarrow D \rightarrow \cdots
\]

If the chain never terminates, then there are no immediate premises.

If there are no immediate premises, then there are no first principles.

If there are no first principles, there is no demonstration.

\[
\boxed{
\text{Demonstration requires immediate premises; therefore regress must stop.}
}
\]

\subsection*{Teaching Point}

Scientific knowledge depends on the termination of explanation in first
principles.

\[
\boxed{
\text{An infinite chain of reasons would prevent completed knowledge.}
}
\]

\section*{29. The Superiority of Universal Demonstration}

Aristotle argues that universal demonstration is superior to particular
demonstration.

This does not mean that Aristotle believes in separated Platonic universals.

Rather, universal demonstration is superior because it gives the cause at the
proper level.

\subsection*{Example}

Particular demonstration:

\[
\text{Isosceles triangles have angles equal to two right angles.}
\]

Universal demonstration:

\[
\text{Triangles have angles equal to two right angles.}
\]

The universal demonstration is better because the attribute belongs to triangle
as such, not to isosceles as such.

\subsection*{Teaching Point}

The best demonstration proves the attribute of the first proper subject.

\[
\boxed{
\text{The right universal gives the right cause.}
}
\]

\section*{30. Universal Does Not Mean Over-General}

Aristotle is not saying that the more general claim is always better.

The correct universal is the level at which the attribute belongs essentially.

\[
\text{Too low: isosceles}
\]

\[
\text{Proper: triangle}
\]

\[
\text{Too high: figure}
\]

Having angles equal to two right angles belongs to triangle as such, not to
isosceles merely, and not to figure generally.

\subsection*{Teaching Point}

Science seeks not the widest possible subject, but the proper explanatory
subject.

\[
\boxed{
\text{The scientific universal is exact, not merely broad.}
}
\]

\section*{31. The Superiority of Affirmative Demonstration}

Aristotle argues that affirmative demonstration is superior to negative
demonstration.

Reasons include:

\begin{itemize}[itemsep=0.25em]
    \item affirmation is prior to negation,
    \item being is prior to not-being,
    \item negative demonstration depends on affirmative structure,
    \item affirmative demonstration is closer to the basic form of proof.
\end{itemize}

\subsection*{Teaching Point}

Affirmative demonstration is more basic because it shows what belongs to what,
not merely what does not belong.

\[
\boxed{
\text{Affirmation is prior to negation in scientific explanation.}
}
\]

\section*{32. Direct Demonstration and Reductio}

Aristotle also ranks direct demonstration above \textit{reductio ad
impossibile}.

A direct demonstration proceeds from premises to conclusion.

\[
P_1 + P_2 \Rightarrow C
\]

A reductio assumes the contradictory and derives impossibility.

\[
\text{not-}C \Rightarrow \bot
\]

\[
\therefore C
\]

\subsection*{Teaching Point}

Reductio may prove, but direct demonstration better displays the prior cause.

\[
\boxed{
\text{Direct explanation is superior to proof through impossibility.}
}
\]

\section*{33. Perception Is Not Demonstration}

Scientific knowledge is not possible through perception alone.

Perception is of:

\[
\text{this}
\]

\[
\text{here}
\]

\[
\text{now}
\]

Science is of:

\[
\text{the universal}
\]

\[
\text{always}
\]

\[
\text{the cause}
\]

\subsection*{Example}

Even if we perceived that a particular triangle has angles equal to two right
angles, we would still seek demonstration.

Why?

Because perception gives the particular, while science grasps the
commensurately universal.

\[
\boxed{
\text{Seeing the fact is not the same as knowing the cause.}
}
\]

\section*{34. The Eclipse Example}

Aristotle says that even if we were on the moon and saw the earth blocking the
sun's light, we would perceive the fact of the eclipse.

But we would not yet possess scientific knowledge unless we grasped the
universal cause.

\[
\text{Perception: the earth is now blocking the sun.}
\]

\[
\text{Science: eclipses occur because light is obstructed in this way.}
\]

\subsection*{Teaching Point}

Perception may start inquiry, but demonstration completes knowledge.

\[
\boxed{
\text{Science requires universal causal explanation beyond perception.}
}
\]

\section*{35. Different Sciences Have Different Basic Truths}

Aristotle argues that not all sciences share the same basic truths.

Each science has principles appropriate to its own subject genus.

\[
\text{Arithmetic} \rightarrow \text{number}
\]

\[
\text{Geometry} \rightarrow \text{magnitude}
\]

\[
\text{Medicine} \rightarrow \text{health and disease}
\]

\subsection*{Teaching Point}

There is no single stock of premises from which every scientific conclusion can
be demonstrated.

\[
\boxed{
\text{Different sciences require different proper principles.}
}
\]

\section*{36. Opinion and Scientific Knowledge}

Scientific knowledge concerns what is necessary and cannot be otherwise.

Opinion concerns what may be otherwise, even if the opinion is true.

\[
\text{Knowledge} = \text{necessary}
\]

\[
\text{Opinion} = \text{can be otherwise}
\]

\subsection*{Teaching Point}

True opinion is not yet scientific knowledge.

\[
\boxed{
\text{To know scientifically is to grasp necessity, not merely truth.}
}
\]

\section*{37. Can One Know and Opine the Same Thing?}

Aristotle says that knowledge and opinion may concern the same subject in a
qualified sense, but not in the same way.

\subsection*{Example}

One may know:

\[
\text{Human is essentially animal.}
\]

This is scientific knowledge if animal is grasped as part of the essence of
human.

One may merely opine:

\[
\text{Human is animal.}
\]

if one does not grasp the essential and necessary connection.

\subsection*{Teaching Point}

The difference lies not only in the object, but in the mode of apprehension.

\[
\boxed{
\text{Knowledge grasps necessity; opinion treats the connection as contingent.}
}
\]

\section*{38. Quick Wit}

Book I ends with quick wit.

Quick wit is the ability to hit upon the middle term instantaneously.

\[
\boxed{
\text{Quick wit} = \text{instant grasp of the middle term}
}
\]

\subsection*{Examples}

A person sees that the moon's bright side is always turned toward the sun and
quickly grasps the cause:

\[
\text{The moon borrows its light from the sun.}
\]

A person sees someone talking with a wealthy man and quickly grasps:

\[
\text{He is borrowing money.}
\]

A person sees two people become friends and quickly grasps:

\[
\text{They share a common enemy.}
\]

\subsection*{Teaching Point}

The whole book has taught us that the middle term is the key to demonstration.
Quick wit is the ability to find that middle term immediately.

\[
\boxed{
\text{To find the middle is to find the cause.}
}
\]

\section*{39. The Architecture of Book I}

The lecture can be summarized as the following movement:

\[
\text{pre-existent knowledge}
\rightarrow
\text{scientific knowledge}
\rightarrow
\text{demonstration}
\rightarrow
\text{first principles}
\rightarrow
\text{essential attributes}
\]

\[
\rightarrow
\text{appropriate sciences}
\rightarrow
\text{fact vs. reasoned fact}
\rightarrow
\text{first figure}
\rightarrow
\text{regress and principles}
\]

\[
\rightarrow
\text{universal demonstration}
\rightarrow
\text{perception, opinion, quick wit}
\]

\subsection*{Teaching Point}

Book I turns formal logic into a theory of science.

\[
\boxed{
\text{The form of inference is not enough; science requires explanatory order.}
}
\]

\section*{40. Summary of Main Ideas}

\begin{enumerate}[itemsep=0.5em]
    \item All teaching and learning by argument begins from pre-existent
    knowledge.

    \item Scientific knowledge is knowledge of the cause as cause.

    \item Demonstration is a syllogism productive of scientific knowledge.

    \item Demonstrative premises must be true, primary, immediate, prior,
    better known, and explanatory.

    \item Not all knowledge is demonstrative; first principles are
    indemonstrable.

    \item Infinite regress and circular demonstration both fail to explain
    scientific knowledge.

    \item Demonstration concerns necessary and essential attributes, not
    accidents as accidents.

    \item The proper scientific universal is the first subject to which the
    attribute belongs as such.

    \item Demonstration must proceed from principles appropriate to the subject
    genus.

    \item Knowledge of the fact differs from knowledge of the reasoned fact.

    \item In scientific demonstration, the middle term is the cause.

    \item The first figure is the most scientific figure.

    \item Universal demonstration is superior to particular demonstration when
    it gives the proper cause.

    \item Affirmative demonstration is superior to negative demonstration, and
    direct demonstration is superior to reductio.

    \item Perception begins inquiry, but scientific knowledge requires
    universal demonstration.

    \item Opinion may be true, but scientific knowledge grasps necessity.

    \item Quick wit is the immediate discovery of the middle term.
\end{enumerate}

\section*{Final Framing}

\[
\boxed{
\textit{Posterior Analytics}, \text{ Book I, turns logic into epistemology.}
}
\]

The \textit{Prior Analytics} teaches us how conclusions follow. The
\textit{Posterior Analytics} teaches us when following becomes knowing.

\[
\boxed{
\text{A syllogism concludes; a demonstration explains.}
}
\]

\end{document}